\documentclass[11pt,a4paper]{article}

\usepackage{kotex}
\usepackage{fontspec}
\usepackage{geometry}
\usepackage{amsmath}
\usepackage{amssymb}
\usepackage{booktabs}
\usepackage{array}
\usepackage{longtable}
\usepackage{tabularx}
\usepackage{graphicx}
\usepackage{caption}
\usepackage{hyperref}
\usepackage{enumitem}
\usepackage{xcolor}
\usepackage{setspace}

\geometry{margin=23mm}
\hypersetup{
  colorlinks=true,
  linkcolor=blue,
  citecolor=blue,
  urlcolor=blue
}
\setlist{nosep}
\emergencystretch=3em
\sloppy

\setmainfont{Noto Serif KR}
\setsansfont{Noto Sans KR}
\setmonofont{Consolas}
\setmainhangulfont{Noto Serif KR}
\setsanshangulfont{Noto Sans KR}
\setmonohangulfont{Noto Sans KR}

\newcommand{\eng}[1]{\texttt{#1}}
\newcommand{\figdir}{figures_full_ko}

\title{\textbf{EmoTrace}\\
\large 생물학적 뇌 구조를 모사한 인공 감성 지능 엔진 개발}
\author{EmoTrace Working Draft v3 Academic}
\date{2026년 5월 5일}

\begin{document}
\onehalfspacing
\maketitle

\begin{abstract}
본 연구는 대규모 언어 모델의 공감 응답이 문장 표면에서는 자연스럽고 안전하지만, 실제 감정 에피소드의 원인, 평가 구조, 원초 정서, 행동 경향을 충분히 보존하지 못하는 문제에서 출발한다. 기존 접근은 감정을 주로 텍스트 라벨, 정적 감정 벡터, 또는 ``공감적으로 답하라''는 문체 지시로 처리한다. 그러나 감정은 이름이나 말투만으로 환원되기 어렵다. 감정은 특정 자극이 내부에서 어떻게 평가되고, 어떤 신경 활성 경로로 전개되며, 어떤 정보가 억제되거나 확산되고, 어떤 행동 경향으로 이어지는지를 포함하는 동적인 상태이다.

EmoTrace는 이러한 문제의식을 바탕으로 감정을 외부 라벨이 아니라 내부 신경 활성 흐름으로 표현하려는 실험적 시스템이다. 연구는 v1부터 v4까지 단계적으로 진행되었다. v1은 입력 문장을 감정/자극 벡터로 변환하고 이를 \eng{EmotionDynamicsNet}에 통과시킨 뒤 시간별 히스토리와 잠재 상태 \eng{z}를 관찰하는 최초의 내부 상태화 파이프라인이었다. v2는 흥분성, 억제성, 조절성 뉴런, trait EMA, memory, cluster, cluster-level rewiring, branch tracking, history encoder, \eng{z -> s} 회귀, prompt generator, style scorer를 포함하는 모듈형 감정 동역학 시스템으로 확장되었다. v3는 dominant branch를 추출하면서 branch collapse와 스타일 목표 편향을 발견하였다. v3.1은 초기 감정 벡터를 감정이 아니라 자극으로 재정의하고, 그 자극이 뉴런망 안에서 만든 trace 자체를 감정 상태 표현 후보로 정의하였다. v4는 trace를 감정 에피소드로 해석하고, targeted episode-sensitive 조건에서 자극만 사용하는 기준선과 비교하였다.

실험 결과, v3 보정 이후 지배 흐름 평균 길이는 1.0539에서 70.4684로 증가했고, 한 칸짜리 흐름 비율은 0.9734에서 0.0948로 감소하였다. v3.1 최종 설정은 흐름 길이 50.475, 한 칸짜리 흐름 비율 0.000, 활성 밀도 0.709412, 감정 구조 분리도 0.238547, 균형 보정 감정 신호 0.136426을 보였으며, 인과 확인 실험의 전체 성공률은 0.916667이었다. v4 targeted 조건에서 \eng{episode\_trace\_v3}는 \eng{stim\_only} 대비 paired n=78, 평균 차이 +1.8308, 95\% 부트스트랩 신뢰구간 [+1.5667, +2.0897], 승/무/패 70/3/5를 보였다. 반면 일반 조건에서는 trace 기반 응답이 항상 우월하지 않았다. 따라서 본 연구의 결론은 EmoTrace가 모든 감정 응답을 해결했다는 것이 아니라, 신경 활성 trace가 감정 상태 표현 후보가 될 수 있으며, 에피소드 정보가 중요한 targeted 조건에서 자극만 사용하는 방식보다 감정 에피소드 충실도를 높일 수 있다는 것이다.
\end{abstract}

\tableofcontents
\newpage

\section{서론}

본 연구는 생성형 언어 모델의 공감하는 응답이 표면적으로는 자연스럽고 안전하지만, 실제 감정의  구조를 충분히 보존하지 못한다는 문제에서 출발한다. EmoTrace는 이 문제를 단순한 문체의 개선에서의 문제가 아니라, 입력 자극과 최종 응답 사이에 계산적으로 추적 가능한 감정 상태 변화 과정의 표현을 나타낼 수 있는가라는 문제로 정의한다.

\subsection{연구 문제}

본 연구의 핵심 질문은 다음과 같다. 첫째, 입력 문장을 감정 라벨로 직접 분류하지 않고, 내부 신경 동역학으로 보내는 파이프라인을 만들 수 있는가. 둘째, 자극이 신경 동역학(뉴런망) 속 에서 지나간 dominant branch 와 시계열 활성 기록을 추출하여 감정 상태를 설명하는 중간 표현으로써 사용할 수 있는가. 셋째, 입력 문장으로 부터 추출한 초기 감정 벡터를 감정 그 자체가 아니라 자극으로 보고, 그 자극이 만든 뉴런망 내의 trace를 감정 상태로써 정의할 수 있는가. 넷째, 이 trace를 감정 에피소드로 해석하면 자극(초기 감정 벡터)만 사용하는 것(\eng{stim\_only})보다 감정이 움직이는 과정에 더 충실한 응답을 만들 수 있는가.

이 질문들은 하나의 성능 지표로 확인될 수 없다. EmoTrace는 단순히 응답을 더 좋게 하는 시스템이 아니라, 감정의 내부 상태가 실제로 생성되었는지, 그 상태가 collapse 없이 유지되는지\textcolor{red}{(죽지 않고 유지 되는지 -붉은 글씨는 뜻을 알아 보기 쉽게 풀이한 것입니다 최종에서는 삭제할 것입니다.)}, 감정축(감정 라벨, 자극) 정보가 남아 있는지, 브랜치 조작에 반응하는지, LLM 응답 생성에 사용했을 때, 어떤 감정 상황에서의 입력에서 실제로 응답 품질이 좋아지는지를 단계적으로 탐구하는 연구이다.

\subsection{주요 기여}

본 연구의 기여는 네 가지로 정리할 수 있다. 첫째, 감정을 출력때의 라벨이 아니라 내부 상태로 보내는 실행 가능한 파이프라인을 구성하였다. 둘째, 흥분성/억제성/조절성 뉴런, trait EMA(특성 지수 이동 평균), memory(기억), clustering(군집화), rewiring, branch tracking을 포함하는 모듈형 감정 동역학 구조를 만들었다. 셋째, EmoTrace v3.1에서 감정의 정의를 새로 하여 초기 감정 벡터를 자극으로, 신경 활성 trace(자극이 뉴런망을 지나간 경로)를 감정 상태의 표현 후보로 제안하였다. 넷째, EmoTrace v4에서 trace를 감정 에피소드로 해석하고 targeted episode-sensitive 조건에서 자극만 사용하는 기준선보다 높은 에피소드 충실도를 보임을 확인하였다.

중요한 점은 본 연구가 EmoTrace가 인간처럼 실제 감정을 느낀다고 주장하지 않는다는 것이다. 본 연구에서 감정 상태란 의식적 주관 경험이 아니라, 입력 자극이 시스템 내부에서 만들어 낸 계산적 활성 구조를 의미한다. 따라서 주장 가능 범위는 ``trace가 감정 상태 표현 후보로 기능할 수 있다''는 계산적 주장에 한정된다.

\begin{figure}[t]
  \centering
  \includegraphics[width=\linewidth]{\figdir/full_research_pipeline.png}
  \caption{EmoTrace 전체 연구 파이프라인과 버전별 역할.}
  \label{fig:academic-full-pipeline}
\end{figure}

\section{선행 연구}

\subsection{기본 감정 이론과 범주형 감정 표현}

감정을 계산적으로 다루는 가장 오래된 방식 중 하나는 감정을 몇 개의 기본 범주로 나누는 접근이다. Ekman은 기본 감정 이론을 통해 일부 감정 표현이 문화권을 넘어 비교적 안정적으로 인식될 수 있다고 주장하였다\cite{ekman1992}. Plutchik은 진화적 관점에서 감정을 생존과 적응에 연결하고, 기본 감정과 그 조합을 설명하는 심리진화론적 모형을 제시하였다\cite{plutchik1980}. 이러한 범주형 접근은 감정을 분류하고 이름 붙이는 데 유용하며, 실제 감정 인식 시스템에서도 직관적인 출력 형식을 제공한다.

그러나 범주형 접근은 본 연구가 다루는 문제를 충분히 설명하지 못한다. 같은 ``분노'' 라벨이라도 공개적 모욕에 대한 분노, 불공정한 평가에 대한 분노, 자기방어적 분노, 타인에게 책임을 돌리는 분노는 서로 다른 평가 구조와 행동 경향을 갖는다. 범주는 감정을 요약하지만, 감정이 어떻게 형성되었는지, 어떤 내부 경로를 거쳤는지, 어떤 행동 경향으로 이어지는지는 설명하지 않는다. EmoTrace는 이 한계를 출발점으로 삼아, 감정 라벨 이전의 내부 활성 과정을 추적하려 한다.

\subsection{차원형 감정 이론과 감정 벡터 표현}

범주형 이론과 달리 Russell의 원형 감정 모형은 감정을 쾌/불쾌와 각성 수준 같은 차원 위에 배치한다\cite{russell1980}. 이 접근은 감정이 서로 완전히 분리된 범주가 아니라 연속적인 공간에서 변할 수 있음을 보여준다. 이후 많은 감정 인식 모델은 감정가, 각성, 지배감 같은 축을 사용하여 감정 상태를 벡터로 표현하였다. 이러한 차원형 접근은 EmoTrace v1/v2의 초기 감정 벡터 설계와도 연결된다.

그러나 정적 감정 벡터 역시 한계를 갖는다. 벡터는 특정 순간의 상태를 요약할 수 있지만, 그 상태가 어떤 자극에서 시작되었고 어떤 내부 동역학을 통해 형성되었는지는 나타내지 않는다. EmoTrace의 v3.1 전환은 바로 이 지점에서 발생하였다. 초기 감정 벡터는 감정 그 자체가 아니라 자극이며, 감정 상태 표현 후보는 그 자극이 뉴런망 안에서 만든 trace라고 재정의하였다.

\subsection{평가 이론과 감정의 과정성}

Scherer의 component process model은 감정을 단일 라벨이나 단일 벡터가 아니라, 사건 평가, 주관적 느낌, 표현, 생리 반응, 동기 및 행동 경향 등이 상호작용하는 과정으로 본다\cite{scherer2009}. Gross의 정서 조절 과정 모형도 감정이 상황, 주의, 평가, 반응의 흐름 속에서 조절될 수 있음을 강조한다\cite{gross1998}. 이러한 과정 중심 관점은 EmoTrace의 trace-as-emotion 가설과 직접적으로 맞닿아 있다.

EmoTrace는 심리학적 평가 과정을 그대로 생물학적으로 재현한다고 주장하지 않는다. 다만 감정이 정적인 이름이나 점수가 아니라, 자극이 내부에서 평가되고 조절되며 행동 경향으로 이어지는 과정이라는 관점을 계산 구조로 옮기고자 한다. v4의 episode 해석에서 사건 원인, 평가 구조, 원초 정서, 행동 경향을 분리한 것도 이러한 선행 이론과 연결된다.

\subsection{정서 컴퓨팅과 기계의 감정 처리}

Picard는 affective computing을 통해 컴퓨터가 인간의 감정을 인식하고, 이해하며, 표현하는 능력이 인간-컴퓨터 상호작용에서 중요하다고 제안하였다\cite{picard1997}. 이후 정서 컴퓨팅 연구는 얼굴 표정, 음성, 생체 신호, 텍스트 등 다양한 입력을 통해 감정을 인식하거나 반응을 생성하는 방향으로 발전하였다.

EmoTrace는 이 전통에 속하지만, 일반적인 감정 인식 모델과 차이가 있다. 많은 정서 컴퓨팅 시스템은 입력에서 감정 라벨 또는 감정 점수를 추정하는 데 초점을 둔다. 반면 EmoTrace는 입력과 출력 사이의 내부 상태를 연구 대상으로 삼는다. 즉, 감정을 인식하는 것에서 그치지 않고, 자극이 내부 뉴런망을 통과하며 어떤 trace를 만들고, 그 trace가 응답 조건으로 어떻게 사용될 수 있는지를 다룬다.

\subsection{공감 대화 모델과 감정 상황 기반 데이터셋}

공감 대화 생성 연구에서는 감정 상황을 기반으로 한 대화 데이터셋과 벤치마크가 제안되었다. Rashkin 등은 \eng{EmpatheticDialogues}를 통해 감정 상황에 기반한 대화 생성 과제를 제시하고, 공감적 응답 생성을 위한 데이터 기반 접근을 발전시켰다\cite{rashkin2019}. 이러한 연구는 감정 라벨뿐 아니라 사용자가 처한 상황 맥락이 공감 응답에 중요하다는 점을 보여준다.

EmoTrace의 문제의식도 이와 유사하게 상황 맥락을 중시한다. 다만 EmoTrace는 감정 상황을 바로 프롬프트나 데이터셋 조건으로 사용하는 데서 멈추지 않는다. 상황에서 자극을 추출하고, 그 자극이 내부 동역학에서 trace를 만들며, trace를 다시 episode로 해석해 응답 조건으로 사용하는 구조를 제안한다. 따라서 EmoTrace는 공감 대화 생성 연구와 연결되지만, 응답 표면이 아니라 내부 상태 표현을 중심에 둔다는 점에서 구분된다.

\subsection{LLM-as-a-Judge 평가와 편향 문제}

최근 LLM 응답 평가는 사람 평가뿐 아니라 LLM을 judge로 사용하는 방식으로 확장되었다. Zheng 등은 MT-Bench와 Chatbot Arena를 통해 LLM-as-a-Judge 방식의 가능성과 한계를 분석하였다\cite{zheng2023}. 이러한 평가는 대규모 비교를 빠르게 수행할 수 있다는 장점이 있지만, judge 모델이 특정 문체, 길이, 자연스러움, 공손함을 선호하는 편향을 가질 수 있다는 문제가 남는다.

EmoTrace에서 이 문제는 특히 중요하다. 본 연구의 목표는 단순히 더 부드럽고 자연스러운 응답을 만드는 것이 아니다. 오히려 기존 LLM 응답이 긍정성, 부드러움, 차분함으로 과도하게 기울어 감정의 거친 핵심을 지우는 문제를 다룬다. 따라서 일반 LLM judge가 선호하는 답변이 EmoTrace의 목표에 부합한다고 곧바로 말할 수 없다. 본 연구가 축 전용 눈가림 judge, targeted episode-fidelity, anti-softening 관점을 도입한 이유가 여기에 있다.

\section{선행 연구로부터 도출되는 문제 제기}

\subsection{선행 연구가 남긴 공백}

앞 절의 선행 연구는 감정을 다루는 여러 중요한 관점을 제공한다. 기본 감정 이론은 감정을 이름 붙일 수 있는 범주로 정리하였고, 차원형 감정 이론은 감정을 연속적인 벡터 공간에서 표현할 수 있게 하였다. 평가 이론과 정서 조절 이론은 감정이 단순한 결과가 아니라 사건 평가와 반응 조절을 포함하는 과정임을 강조하였다. 정서 컴퓨팅과 공감 대화 연구는 기계가 감정을 인식하고 공감적 응답을 생성하는 문제를 본격적인 계산 과제로 만들었다. 또한 LLM-as-a-Judge 연구는 대규모 응답 평가를 가능하게 하였지만, 평가자가 특정 문체와 자연스러움을 선호할 수 있다는 한계도 드러냈다.

이러한 선행 연구를 종합하면 하나의 공백이 드러난다. 기존 연구는 감정을 라벨로 분류하거나, 벡터로 요약하거나, 상황 기반 응답을 생성하는 데 강점을 보인다. 그러나 입력 자극이 시스템 내부에서 어떤 시간적 활성 경로를 만들고, 그 경로가 감정축 정보와 어떻게 연결되며, 그 내부 경로가 응답 조건으로 어떻게 사용될 수 있는지를 충분히 다루지는 못한다. 즉, 감정의 이름, 감정의 좌표, 감정 상황, 공감 응답은 연구되어 왔지만, ``자극이 내부 동역학 속에서 감정 상태 표현으로 형성되는 과정''은 상대적으로 비어 있다.

\subsection{LLM 공감 응답의 표면적 자연스러움과 구조적 결핍}

이 공백은 LLM 공감 응답에서 특히 뚜렷하게 나타난다. 생성형 언어 모델은 사용자의 감정적 발화를 인식하고 비교적 자연스러운 위로, 조언, 공감 응답을 생성할 수 있다. 그러나 이러한 응답은 종종 표면적으로는 친절하고 안전하지만, 사용자가 실제로 겪는 감정 상태의 핵심을 놓친다. 특히 공개적 모욕, 부당한 평가, 자기비난, 타인 책임, 사회적 위협, 행동 보류와 같은 상황에서는 감정의 본질이 단순한 긍정/부정이나 슬픔/분노 라벨에 있지 않다. 중요한 것은 왜 그 감정이 생겼는지, 어떤 평가 구조를 거쳤는지, 어떤 원초 정서가 남아 있는지, 어떤 행동 경향이 억제되거나 보류되고 있는지이다.

기존 LLM의 공감 응답은 대체로 부드러움, 차분함, 긍정성, 협조성을 높이는 방향으로 생성된다. 이 특성은 일반적인 안전성 관점에서는 유용하지만, 감정 에피소드의 충실도 관점에서는 문제가 될 수 있다. 예를 들어 사용자가 공개적으로 모욕당한 상황을 말했을 때, 모델이 곧바로 ``상대도 사정이 있었을 수 있다''거나 ``너무 마음에 담아두지 말라''고 답한다면, 문장은 안전할 수 있으나 감정의 구조는 손상된다. 이 경우 핵심은 단순한 위로가 아니라 공개적 평가 위협, 수치심, 분노, 항의 욕구, 관계 위험 때문에 행동을 보류하는 상태이다.

따라서 본 연구는 공감 응답의 문제를 ``더 따뜻한 문장을 생성하는 문제''로 보지 않는다. 오히려 지나치게 따뜻하고 온건한 문장이 감정의 거친 핵심을 지울 수 있다는 점을 문제로 삼는다. 감정 응답에서 필요한 것은 무조건적인 부드러움이 아니라, 감정의 원인과 평가 구조를 보존한 뒤 안전하게 표면화하는 것이다.

\subsection{라벨과 벡터를 넘어 내부 trace가 필요한 이유}

감정 라벨은 결과 요약으로 유용하지만, 감정이 만들어지는 내부 과정을 설명하지 못한다. 감정 벡터는 라벨보다 세밀하지만, 여전히 특정 순간의 요약값에 가깝다. 벡터는 감정 성분의 크기를 나타낼 수 있으나, 그 성분이 어떤 경로로 형성되었고 어떤 정보가 중간에서 억제되었는지는 보여주지 않는다. 이 한계는 선행 연구의 약점이라기보다, 라벨과 벡터가 원래 담당하는 역할의 한계이다.

EmoTrace는 이 지점에서 문제를 다시 정의한다. 초기 감정 벡터는 감정 그 자체가 아니라 입력 문장이 시스템에 가하는 자극이다. 감정 상태 표현 후보는 그 자극이 뉴런망 안에서 만든 시간적 활성 기록, 즉 trace이다. 이 관점에서는 감정을 잘 분류하는 것보다, 자극이 내부에서 어떤 경로를 만들고 그 경로가 어떤 감정축 정보를 보존하는지를 확인하는 것이 중요해진다.

\subsection{내부 상태 모델링의 필요성}

내부 상태 모델링은 단순한 복잡화가 아니라 실패를 분해하기 위한 장치이다. 응답이 어색할 때 그 원인은 입력 자극 추출, 뉴런망 동역학, branch 기록, history encoding, \eng{z -> s} 변환, prompt generation, LLM judge 편향 중 어디에나 있을 수 있다. 입력과 출력만 보면 이 원인을 분리할 수 없다.

EmoTrace는 입력과 출력 사이에 자극, 동역학, branch, trace, episode, 응답 조건화, 평가라는 중간 단계를 둔다. 이 구조는 성능 향상만이 아니라 진단 가능성을 제공한다. v3에서 branch collapse를 발견할 수 있었던 것도 내부 경로를 기록하는 구조가 있었기 때문이다. 따라서 EmoTrace의 연구 대상은 최종 응답만이 아니라 그 응답을 가능하게 하는 내부 활성 과정 전체이다.

\section{개념적 정의와 연구 가설}

\subsection{자극, 상태, trace의 구분}

v3.1 이후 EmoTrace에서 가장 중요한 정의는 자극과 감정 상태의 분리이다. 초기 감정 벡터는 더 이상 감정 그 자체가 아니다. 그것은 입력 문장이 시스템에 가하는 최초 압력, 즉 자극이다. 감정 상태 표현 후보는 그 자극이 뉴런망 안에서 만든 시간적 활성 기록, 즉 trace이다.

이 구분은 연구의 논리 구조를 바꾼다. 초기 벡터를 감정으로 보면 연구의 목표는 더 좋은 감정 벡터를 만드는 일이 된다. 반대로 초기 벡터를 자극으로 보면 연구의 목표는 그 자극이 내부에서 어떤 경로를 만들고, 그 경로가 어떤 구조를 유지하며, 그 구조가 감정축 정보와 어떻게 연결되는지를 분석하는 일이 된다.

\subsection{trace-as-emotion 가설}

본 연구의 중심 가설은 다음과 같다.

\begin{quote}
입력 문장에서 추출된 초기 감정 벡터는 감정 자체가 아니라 자극이며, 감정 상태 표현 후보는 그 자극이 뉴런망 안에서 만든 신경 활성 trace이다.
\end{quote}

이 가설은 인간의 주관적 감정을 기계가 동일하게 경험한다는 주장이 아니다. 여기서 감정 상태는 계산적으로 관찰 가능한 내부 활성 구조이다. 따라서 trace-as-emotion 가설은 감정의 존재론적 주장이라기보다, 감정 응답 생성을 위한 중간 표현 가설이다.

\subsection{에피소드 조건화 가설}

v4에서는 trace를 감정 에피소드로 해석하였다. 에피소드란 감정 이름 하나가 아니라, 자극의 사건 원인, 평가 구조, 원초 정서, 행동 경향, 표현 시 위험을 포함하는 중간 표현이다. 에피소드 조건화 가설은 다음과 같다.

\begin{quote}
감정 원인과 행동 경향이 중요한 입력에서는, trace에서 해석한 episode 정보를 응답 조건으로 사용할 때 자극만 사용하는 기준선보다 감정 에피소드 충실도가 높아질 수 있다.
\end{quote}

이 가설은 모든 일반 입력에서 trace 기반 응답이 우수하다는 뜻이 아니다. 오히려 본 연구는 일반 조건과 targeted 조건을 구분한다. trace 정보가 중요하지 않은 일반 입력에서는 큰 차이가 나지 않을 수 있다. 반대로 공개적 모욕, 자기비난, 타인 책임, 사회적 위협, 행동 보류처럼 에피소드 구조가 중요한 입력에서는 trace 정보가 유용할 가능성이 크다.

\begin{figure}[t]
  \centering
  \includegraphics[width=0.92\linewidth]{\figdir/emotion_definition_shift.png}
  \caption{감정 정의의 변화. v1/v2는 결과 벡터 중심 해석을 사용했고, v3는 대표 경로의 중요성을 발견했으며, v3.1은 trace 자체를 감정 상태 표현 후보로 재정의했다.}
  \label{fig:academic-definition-shift}
\end{figure}

\section{시스템 발전 과정}

\begin{figure}[t]
  \centering
  \includegraphics[width=\linewidth]{\figdir/version_timeline.png}
  \caption{EmoTrace 연구 버전의 개념 발전.}
  \label{fig:academic-timeline}
\end{figure}

\subsection{v1: 내부 상태화 파이프라인의 형성}

v1은 감정을 텍스트 라벨로만 예측하지 않고, 입력 문장에서 얻은 정서 자극을 내부 동역학으로 통과시킨 뒤 새로운 잠재 상태로 바꾸려는 최초의 실행 가능한 구조였다. 구조는 입력 문장, 감정/자극 벡터, \eng{EmotionDynamicsNet}, 시간별 히스토리, GRU history encoder, 잠재 상태 \eng{z}, GUI 관찰로 이어졌다.

\begin{figure}[t]
  \centering
  \includegraphics[width=\linewidth]{\figdir/v1_pipeline_gui.png}
  \caption{v1 파이프라인. 입력 문장은 감정/자극 벡터로 변환되고, \eng{EmotionDynamicsNet}의 시간별 히스토리를 거쳐 잠재 상태 \eng{z}로 압축된다.}
  \label{fig:academic-v1}
\end{figure}

v1의 의미는 단순한 프로토타입에 있지 않다. v1은 LLM 공감 문제를 ``좋은 문장을 쓰는 문제''에서 ``내부 상태를 생성하는 문제''로 옮겼다. 이때 \eng{z}는 단순한 분류 결과가 아니라, 입력 감정이 뉴런망 안에서 변형된 뒤 남은 내부 상태의 압축 표현으로 해석되었다.

v1 GUI는 연구 관찰 장치였다. 연구자는 입력, 추론 결과, \eng{z}, 히스토리 변화를 한 화면에서 확인할 수 있었다. 이 관찰 경험은 이후 trace와 branch를 더 자세히 기록해야 한다는 필요성으로 이어졌다. 따라서 v1의 성과는 정량 성능이 아니라, 감정을 내부 상태로 보낼 수 있는 실행 가능한 파이프라인과 관찰 도구를 만들었다는 점이다.

동시에 v1은 결과 벡터 중심 해석의 한계를 남겼다. 히스토리와 \eng{z}는 관찰할 수 있었지만, 어떤 뉴런이 어떤 경로로 활성화되었고 어떤 정보가 유지되거나 사라졌는지는 충분히 기록되지 않았다. 이 한계가 v2의 모듈화와 v3의 branch tracking을 요구했다.

\subsection{v2: 모듈형 감정 동역학 시스템}

v2는 v1의 단일 흐름을 분해 가능한 모듈형 구조로 확장하였다. v2에는 텍스트에서 4차원 제어 벡터를 만드는 encoder, 흥분성/억제성/조절성 뉴런을 포함한 dynamics, trait EMA, memory, cluster, cluster-level rewiring, branch tracking, history encoder, \eng{z -> s} style regression, prompt generator, style scorer가 포함되었다.

\begin{figure}[t]
  \centering
  \includegraphics[width=\linewidth]{\figdir/v2_module_structure.png}
  \caption{v2 모듈 구조. 모듈 분리는 이후 v3/v3.1/v4의 진단 실험을 가능하게 했다.}
  \label{fig:academic-v2-modules}
\end{figure}

v2의 독창성은 감정 벡터를 단순 회귀 모델이나 MLP에 넣은 것이 아니라, 기능이 다른 뉴런 집단에 통과시켰다는 점이다. 흥분성 뉴런은 감정 신호를 확산시키고, 억제성 뉴런은 과도한 활성화를 막으며, 조절성 뉴런은 흐름의 민감도와 방향을 조정한다. trait EMA는 한 번의 입력만이 아니라 이전 상태의 흔적이 다음 반응에 영향을 주도록 만들었고, memory는 내부 상태가 순간 입력에만 의존하지 않도록 하였다. clustering과 rewiring은 고정된 선형 경로가 아니라 상황에 따라 다른 경로가 형성될 수 있도록 설계되었다.

\begin{figure}[t]
  \centering
  \includegraphics[width=0.9\linewidth]{\figdir/v2_neuron_types.png}
  \caption{v2 뉴런 타입 구조. 흥분성, 억제성, 조절성 뉴런은 각각 신호 확산, 과활성 억제, 흐름 조절의 역할을 맡는다.}
  \label{fig:academic-v2-neurons}
\end{figure}

v2의 가장 중요한 성과는 최종 성능 수치가 아니라 문제를 나누어 볼 수 있는 구조를 만든 것이다. v1에서는 결과가 나빠졌을 때 원인이 encoder인지, dynamics인지, memory인지, 응답 생성 조건인지 분리하기 어려웠다. v2에서는 encoder, dynamics, branching, history encoder, tone regressor, prompt generator, style scorer가 분리되면서 각 모듈을 따로 진단할 수 있게 되었다.

또한 v2의 branch tracking은 trace-as-emotion 관점의 씨앗이었다. 당시 branch는 감정 그 자체로 정의되지 않았고, 내부 경로를 기록하는 보조 장치에 가까웠다. 그러나 이 장치가 있었기 때문에 v3에서 dominant branch를 추출할 수 있었고, branch collapse라는 결정적 병목도 발견할 수 있었다.

\subsection{v3: dominant branch와 collapse의 발견}

v3는 자극이 뉴런망 안에서 지나간 대표 경로를 추출하기 시작한 단계이다. 입력 문장은 자극 벡터로 변환되고, 클러스터 기반 뉴런망 동역학을 거치며, tick별 활성과 firing edge가 기록된다. 이후 이 기록에서 dominant branch를 추출하고, 지배 흐름은 잠재 표현 \eng{z}로 압축되며, \eng{z}는 응답 스타일 표현 \eng{s}로 회귀된다.

v3의 핵심 발견은 branch collapse였다. 초기 full export 51,628개 기준 지배 흐름 평균 길이는 1.0539였고, 한 칸짜리 흐름 비율은 0.9734였다. 이는 대부분의 흐름이 한 tick 또는 한 노드 수준에서 종료되었음을 의미한다. 이러한 상태에서는 branch를 감정 상태 표현의 근거로 사용할 수 없다.

\begin{figure}[t]
  \centering
  \includegraphics[width=0.92\linewidth]{\figdir/branch_collapse_recovery.png}
  \caption{v3 보정 전후의 지배 흐름 붕괴 완화.}
  \label{fig:academic-branch-recovery}
\end{figure}

branch collapse는 단순한 구현 오류가 아니라 개념적 위기였다. EmoTrace가 내부 경로를 감정 상태의 근거로 삼으려면, 그 경로가 실제로 존재해야 한다. 지배 흐름이 대부분 한 칸에서 끝난다면 내부 감정 전개가 거의 없는 셈이다. 따라서 v3의 보정은 성능 튜닝이 아니라 연구 가설을 유지하기 위한 필수 단계였다.

보정 이후 지배 흐름 평균 길이는 70.4684로 증가했고, 한 칸짜리 흐름 비율은 0.0948로 낮아졌다. 이는 trace 분석을 위한 최소 동역학 조건이 회복되었음을 의미한다. 이 과정에서 평균 흐름 길이, 한 칸짜리 흐름 비율, 최대 길이, 95백분위 길이, 첫 활성 tick, 활성 창 길이와 같은 동역학 지표가 연구의 핵심 평가 도구가 되었다.

v3는 또한 스타일 목표 편향을 발견했다. 스타일 목표가 긍정성, 부드러움, 차분함, 협조성으로 과도하게 치우쳐 있었고, 적대성, 원망, 절망, 불안정성, 두려움, 수치심은 거의 0에 가까웠다. 일반 챗봇 기준에서는 좋아 보일 수 있으나, EmoTrace의 목표와는 충돌한다. 감정 에피소드의 핵심이 분노, 수치심, 두려움, 원망에 있는데 이를 모두 부드럽게 덮으면 감정 상태의 핵심이 사라진다.

\begin{figure}[t]
  \centering
  \includegraphics[width=0.92\linewidth]{\figdir/style_bias_korean.png}
  \caption{v3 스타일 목표 편향. 온건한 축은 높고 날것의 감정축은 낮게 나타났다.}
  \label{fig:academic-style-bias}
\end{figure}

\subsection{v3.1: trace-as-emotion 정의 전환}

v3.1은 EmoTrace의 중심 개념이 바뀐 단계이다. v1/v2에서는 감정을 결과 벡터 또는 잠재 표현 \eng{z}로 보는 경향이 강했다. v3에서는 대표 경로가 중요해졌지만, 아직 trace 자체가 감정으로 정의되지는 않았다. v3.1에서는 최초 감정 벡터를 감정이 아니라 자극으로 보고, 그 자극이 뉴런망 안에서 만든 trace 자체를 감정 상태 표현 후보로 정의하였다.

이 전환은 용어 변경이 아니라 연구 대상의 변경이다. 결과 벡터 중심 관점에서는 \eng{z}가 얼마나 좋은지가 핵심이다. trace 중심 관점에서는 \eng{z}가 만들어지는 동안 어떤 뉴런이 켜지고, 어떤 경로가 유지되며, 어떤 정보가 억제되거나 확산되는지가 핵심이다. 따라서 v3.1은 생성 품질보다 표현 검증을 우선하였다.

\begin{figure}[t]
  \centering
  \includegraphics[width=0.9\linewidth]{\figdir/v31_adaptive_metrics.png}
  \caption{v3.1 적응형 설정의 주요 동역학 및 표현 지표.}
  \label{fig:academic-v31-metrics}
\end{figure}

v3.1은 trace가 감정 상태 표현 후보가 되기 위한 조건을 지표로 정의하였다. trace는 충분한 흐름 길이를 가져야 하고, 한 칸짜리 흐름 비율이 낮아야 하며, 활성 밀도가 너무 낮거나 높지 않아야 한다. 또한 감정 구조 분리도와 균형 보정 감정 신호를 통해 감정축 정보가 남아 있는지 확인해야 한다.

\begin{figure}[t]
  \centering
  \includegraphics[width=0.92\linewidth]{\figdir/v31_causal_results.png}
  \caption{v3.1 trace 인과 조작 실험 결과.}
  \label{fig:academic-v31-causal}
\end{figure}

인과 조작 실험은 trace가 단순 로그가 아니라 조작 가능한 중간 표현인지 확인하기 위한 실험이었다. 방향 교란은 특정 감정축의 방향을 바꾸고, 정보 중립화는 특정 축 단서를 neutral로 바꾼다. 축 전용 눈가림 judge는 일반 품질, 유창성, 친절함, 자연스러움, 공손함 같은 단서를 보지 못하게 하고, 목표 감정축의 변화만 판정하도록 설계되었다. 동일 응답 무효 비교가 40/40 tie였다는 점은 judge가 억지로 차이를 만들어 낸 것이 아님을 보조적으로 보여준다.

\subsection{v4: 에피소드 해석과 targeted 평가}

v4는 trace를 감정 에피소드로 해석하고, 응답 생성 조건으로 사용하는 단계이다. 에피소드 해석은 감정 이름 하나를 출력하는 것이 아니라, 자극의 사건 원인, 평가 구조, 원초 정서, 행동 경향, 표현상의 위험을 중간 표현으로 구성한다. 예를 들어 공개적 모욕 상황에서 중요한 것은 단순한 ``분노'' 라벨이 아니라 공개성, 평가 위협, 수치심, 항의 욕구, 관계 위험 때문에 행동을 보류하는 구조이다.

\begin{figure}[t]
  \centering
  \includegraphics[width=0.95\linewidth]{\figdir/episode_conditioning_workflow.png}
  \caption{v4 episode conditioning 흐름. trace는 응답 가능한 사건 구조로 번역된다.}
  \label{fig:academic-episode-workflow}
\end{figure}

v4에서는 일반 조건과 targeted episode-sensitive 조건을 구분하였다. 일반 조건에서는 trace 기반 응답이 항상 우수하지 않았다. 이는 EmoTrace의 실패가 아니라 claim을 좁히는 근거이다. trace의 강점은 감정 원인, 평가 구조, 행동 경향을 보존하는 데 있다. 따라서 이 정보가 중요하지 않은 일반 입력에서 큰 차이가 나지 않는 것은 자연스럽다. 반대로 공개적 모욕, 자기비난, 타인 책임, 사회적 위협, 행동 보류처럼 에피소드 구조가 중요한 입력에서는 trace 정보가 실제로 도움이 될 수 있다.

\section{실험 설계}

\subsection{실험 흐름}

본 연구의 실험은 단일 벤치마크 성능 평가가 아니라, 연구 가설이 변화하는 과정에 따라 누적적으로 설계되었다. v1/v2에서는 내부 상태화 파이프라인과 모듈형 구조의 성립을 확인하였다. v3에서는 branch collapse를 발견하고 dynamics를 보정하였다. v3.1에서는 trace가 감정 상태 표현 후보로 기능할 수 있는지 동역학, 표현, 인과 지표로 검증하였다. v4에서는 trace를 episode로 해석하고 응답 조건으로 사용했을 때 targeted 조건에서 효과가 있는지 비교하였다.

\begin{figure}[t]
  \centering
  \includegraphics[width=\linewidth]{\figdir/experiment_before_after_map.png}
  \caption{주요 실험의 이전/이후 구조.}
  \label{fig:academic-before-after}
\end{figure}

\subsection{초기 구조 실험}

초기 구조 실험의 목적은 정량 성능을 극대화하는 것이 아니라, 감정을 내부 상태로 보낼 수 있는지 확인하는 것이었다. v1은 입력 문장을 감정/자극 벡터로 바꾸고, 이를 동역학 네트워크에 통과시키며, 시간별 히스토리와 \eng{z}를 관찰하였다. v2는 이를 모듈화하여 이후 실험에서 각 부분을 분리 진단할 수 있도록 만들었다.

이 단계의 성과는 ``실행 가능한 구조''이다. v1/v2가 없었다면 v3에서 branch collapse를 발견할 수 없었고, v3.1에서 trace를 감정 상태 표현 후보로 검증할 수도 없었다. 따라서 v1/v2의 의미는 높은 성능 수치가 아니라 연구 가능한 내부 관찰 구조를 만든 데 있다.

\subsection{branch/dynamics 보정 실험}

branch/dynamics 보정 실험은 v3에서 발견된 branch collapse를 해결하기 위한 실험이다. 초기 지배 흐름 평균 길이 1.0539와 한 칸짜리 흐름 비율 0.9734는 trace 기반 감정 표현 가설을 위협하였다. 보정 이후 평균 길이 70.4684, 한 칸짜리 흐름 비율 0.0948을 얻으면서 trace 분석의 최소 조건을 회복하였다.

이 실험은 성공적인 응답 생성을 직접 증명하지 않는다. 그러나 trace를 감정 상태 표현 후보로 다루기 위한 전제 조건을 검증한다. 내부 경로가 충분히 길고 유지되지 않는다면, 이후 표현 검증과 인과 조작은 의미를 갖기 어렵다.

\subsection{스타일 편향 분석}

스타일 편향 분석은 v3에서 나타난 응답 표면화의 문제를 다룬다. 스타일 목표가 부드러움, 차분함, 협조성, 긍정성에 치우치고, 적대성, 원망, 절망, 불안정성, 두려움, 수치심은 거의 0에 가까웠다. 이는 일반적인 LLM judge가 선호하는 자연스러운 답변과 EmoTrace의 목표가 다를 수 있음을 보여준다.

본 연구는 이 문제를 단순히 ``나쁜 편향''으로만 보지 않았다. 이 편향은 평가 기준과 생성 기준을 분리해야 한다는 사실을 드러냈다. LLM judge가 부드러운 답변을 선호한다면, 감정 에피소드 충실도를 평가하기 위한 별도의 targeted 지표와 축 전용 눈가림 judge가 필요하다.

\subsection{trace 표현 검증}

trace 표현 검증은 v3.1의 핵심 실험이다. 최종 적응형 설정은 흐름 길이 50.475, 한 칸짜리 흐름 비율 0.000, 활성 밀도 0.709412, 감정 구조 분리도 0.238547, 균형 보정 감정 신호 0.136426을 보였다. 이 값들은 trace가 단순히 길어진 로그가 아니라, 감정축 구조를 일부 보존하는 활성 표현일 가능성을 보여준다.

여기서 중요한 점은 하나의 지표만으로 결론을 내리지 않았다는 것이다. 흐름이 길기만 하고 활성 밀도가 지나치게 낮으면 감정 상태로 보기 어렵다. 활성 밀도가 높기만 하고 구조 분리도가 없으면 모든 자극이 비슷하게 퍼지는 것일 수 있다. 균형 보정 신호가 없으면 데이터 불균형 때문에 특정 축이 과장될 수 있다. 따라서 v3.1은 동역학 지표와 표현 지표를 함께 사용하였다.

\subsection{trace 인과 조작 실험}

인과 조작 실험은 trace가 감정축 해석에 실제로 영향을 주는지 확인하기 위한 실험이다. 방향 교란은 특정 축의 방향을 의도적으로 바꾸고, 강한 정보 중립화는 특정 축 정보를 neutral로 만든다. 그 뒤 축 전용 눈가림 judge가 목표 축 변화만을 판정한다.

전체 성공률 0.916667은 trace가 조작에 반응한다는 초기 증거이다. 다만 이 결과만으로 trace가 완전한 감정 상태 표현이라고 주장할 수는 없다. judge 편향, 데이터 규모, 축 설계의 한계가 남아 있다. 따라서 본 연구는 이를 ``초기 인과 증거''로 해석한다.

\subsection{episode 응답 생성 및 targeted 비교}

v4에서는 trace를 episode로 해석한 뒤 응답 조건으로 사용하였다. 비교 조건은 자극만 사용하는 \eng{stim\_only}, 일반 episode trace, 보정된 \eng{episode\_trace\_v3} 등이었다. 일반 조건에서는 trace 기반 응답이 항상 우월하지 않았지만, targeted episode-sensitive set에서는 \eng{episode\_trace\_v3}가 \eng{stim\_only}보다 뚜렷하게 우세했다.

\begin{figure}[t]
  \centering
  \includegraphics[width=0.9\linewidth]{\figdir/v4_targeted_scores.png}
  \caption{v4 targeted episode-sensitive 입력에서의 조건별 점수.}
  \label{fig:academic-targeted-scores}
\end{figure}

\begin{figure}[t]
  \centering
  \includegraphics[width=0.9\linewidth]{\figdir/v4_paired_win_tie_loss.png}
  \caption{v4 쌍대 우월성 평가.}
  \label{fig:academic-paired}
\end{figure}

\section{평가 지표}

\begin{figure}[t]
  \centering
  \includegraphics[width=\linewidth]{\figdir/metric_taxonomy.png}
  \caption{EmoTrace 지표 체계. trace가 생성되었는지, 의미가 있는지, 응답에 도움이 되는지를 분리하여 평가한다.}
  \label{fig:academic-metrics}
\end{figure}

\subsection{동역학 지표}

동역학 지표는 trace가 감정 상태 표현 후보로 사용할 수 있을 만큼 살아 있는지를 판단한다. 주요 지표는 감정 흐름 길이, 한 칸짜리 흐름 비율, 활성 밀도, 최대 tick 도달 비율, 첫 활성 tick, 활성 창 길이이다. 이 지표들은 branch collapse를 발견하고 완화하는 과정에서 핵심적인 역할을 했다.

\subsection{표현 지표}

표현 지표는 trace가 감정축 구조를 담고 있는지 평가한다. v3.1에서는 감정가, 사회적 방향, 행동 경향, 평가 계열, 통제감 상태와 같은 축을 사용하였다. 단순 정확도보다 중요한 것은 데이터 불균형을 고려한 균형 보정 신호이다. 특정 축이 데이터 분포 때문에 우세해지는 것을 막기 위해 균형 보정 감정 신호를 사용하였다.

\subsection{인과 지표}

인과 지표는 trace 조작이 감정축 해석을 움직이는지 평가한다. 방향 교란, 정보 중립화, 동일 응답 무효 비교, 축 전용 눈가림 judge가 포함된다. 이 지표들은 trace가 단순한 기록인지, 아니면 해석 가능한 중간 표현인지 구분하기 위해 도입되었다.

\subsection{생성 지표}

생성 지표는 일반 응답 품질과 targeted episode-fidelity로 나뉜다. 일반 응답 품질은 유창성, 자연스러움, 공손함, 전체 품질을 포함할 수 있다. 그러나 EmoTrace의 핵심 지표는 targeted episode-fidelity이다. 여기에는 상황 평가 충실도, 원초 정서 반영, 과도한 온건화 방지, 행동 경향 적합도, 구체성, 자연스러움, 전체 선호가 포함된다.

\section{결과}

\subsection{내부 상태화 구조의 성립}

v1/v2는 최종 성능 우월성을 입증한 단계가 아니라, 내부 상태화 구조가 실행 가능함을 보인 단계이다. v1은 입력 문장을 자극 벡터로 바꾸고 이를 동역학 네트워크에 통과시켜 히스토리와 \eng{z}를 관찰하였다. v2는 이를 모듈형 구조로 확장하여 이후 실험에서 원인을 분리할 수 있게 하였다.

이 결과는 정량 점수보다 구조적 의미가 크다. 감정을 내부 상태로 보낼 수 없었다면 trace를 논할 수 없고, branch tracking이 없었다면 collapse를 발견할 수 없으며, 모듈 분리가 없었다면 어떤 부분을 수정해야 하는지도 알 수 없었을 것이다.

\subsection{branch collapse 완화}

v3 보정 이후 지배 흐름 평균 길이는 1.0539에서 70.4684로 증가했고, 한 칸짜리 흐름 비율은 0.9734에서 0.0948로 감소하였다. 이 변화는 trace가 분석 가능한 길이와 구조를 갖게 되었음을 의미한다. 따라서 v3 보정은 v3.1 trace 표현 검증의 전제 조건이었다.

\subsection{스타일 편향 확인}

v3 스타일 목표 분석에서는 부드러움, 차분함, 협조성, 긍정성이 높고, 적대성, 원망, 절망, 불안정성, 두려움, 수치심이 낮은 편향이 확인되었다. 이는 LLM judge가 선호하는 자연스럽고 안전한 문장이 EmoTrace의 목표와 항상 일치하지 않음을 보여준다. 이 결과는 이후 targeted 지표와 anti-softening 관점을 도입하는 이유가 되었다.

\subsection{trace 표현 검증}

v3.1의 최종 적응형 설정은 흐름 길이 50.475, 한 칸짜리 흐름 비율 0.000, 활성 밀도 0.709412, 감정 구조 분리도 0.238547, 균형 보정 감정 신호 0.136426을 보였다. 이는 trace가 단순 로그가 아니라 감정축 구조를 일부 담고 있음을 시사한다.

\subsection{인과 조작 결과}

방향 교란과 강한 정보 중립화 실험은 trace 조작이 감정축 해석에 영향을 줄 수 있음을 보였다. 전체 성공률은 0.916667이었다. 동일 응답 무효 비교에서 40/40 tie가 나온 점은 judge가 무조건 차이를 만들어내는 방식으로 작동하지 않았음을 보여준다.

\subsection{일반 조건 실패와 targeted 조건 성공}

v4 일반 조건에서는 \eng{episode\_trace}가 \eng{stim\_only}보다 우세하지 않았다. \eng{episode\_trace} vs \eng{stim\_only}의 평균 차이는 -0.081416이었고 신뢰구간은 0을 포함했다. 반면 targeted 조건에서는 \eng{episode\_trace\_v3}가 \eng{stim\_only}보다 평균 차이 +1.8308, 95\% 부트스트랩 신뢰구간 [+1.5667, +2.0897], 승/무/패 70/3/5로 우세했다.

\begin{figure}[t]
  \centering
  \includegraphics[width=0.86\linewidth]{\figdir/general_vs_targeted_after.png}
  \caption{일반 조건과 targeted 조건의 이전/이후 비교.}
  \label{fig:academic-general-targeted}
\end{figure}

\begin{figure}[t]
  \centering
  \includegraphics[width=0.88\linewidth]{\figdir/targeted_metric_deltas_after.png}
  \caption{targeted 조건에서 지표별 향상량.}
  \label{fig:academic-targeted-deltas}
\end{figure}

이 대비는 본 연구의 claim을 정교하게 만든다. EmoTrace는 일반 감정 응답에서 항상 우수한 만능 생성기가 아니다. 그러나 감정 원인, 평가 구조, 행동 경향이 중요한 targeted 조건에서는 trace 기반 episode 정보가 자극만 사용하는 방식보다 더 높은 에피소드 충실도를 제공할 수 있다.

\section{논의}

\subsection{결과 벡터 중심 감정관에서 과정 중심 감정관으로}

EmoTrace의 가장 중요한 변화는 감정을 결과 벡터로 보는 관점에서 과정으로 보는 관점으로 이동한 것이다. v1/v2에서는 감정 벡터가 뉴런망을 통과한 결과 상태 또는 잠재 표현 \eng{z}가 감정 상태로 해석되었다. 이 관점은 초기 구현에는 유용했지만, 감정이 어떤 경로로 형성되었는지 설명하지 못했다.

v3.1 이후 감정은 결과값이 아니라 trace로 재정의되었다. 즉, 감정 상태 표현 후보는 최종 벡터 하나가 아니라, 자극이 뉴런망 안에서 어떤 뉴런을 켜고, 어떤 경로를 유지하며, 어떤 정보를 억제하거나 확산시키는지의 과정이다. 이 전환이 EmoTrace를 단순 감정 분류기나 단순 프롬프트 템플릿과 구분한다.

\subsection{실패가 연구를 조직한 방식}

본 연구는 처음부터 완성된 구조로 출발하지 않았다. v1은 내부 상태화의 가능성을 보였지만 세밀한 경로 분석이 부족했다. v2는 모듈화를 이루었지만 branch를 감정 본체로 정의하지는 못했다. v3는 dominant branch를 추출했지만 branch collapse를 발견했다. v3 스타일 분석은 과도한 온건화 편향을 드러냈다. v4 일반 조건에서는 trace 기반 응답이 항상 우수하지 않았다.

이 실패들은 연구의 약점이 아니라 다음 단계의 근거가 되었다. branch collapse 때문에 동역학 지표가 생겼고, 스타일 편향 때문에 targeted 지표와 anti-softening 관점이 생겼으며, 일반 조건 실패 때문에 claim을 targeted episode-sensitive 조건으로 좁혔다. 좋은 논문은 성공만 나열하는 것이 아니라, 실패가 어떻게 정의를 바꾸고 실험을 정교화했는지를 보여야 한다.

\subsection{LLM judge 편향에 대한 해석}

LLM judge는 부드럽고 자연스러운 답변을 선호할 수 있다. 따라서 LLM judge만으로 EmoTrace의 목표 달성 여부를 판단하면, 오히려 기존 LLM 스타일에 가까운 답변이 좋은 답변으로 평가될 위험이 있다. 본 연구는 이 문제를 인식하고 축 전용 눈가림 judge와 targeted episode-fidelity 지표를 사용하였다.

그러나 LLM judge 편향이 완전히 제거된 것은 아니다. 향후에는 사람 평가, 다중 judge 비교, blind pairwise evaluation, 감정축별 독립 평가가 필요하다. 현재 결과는 LLM judge 기반의 초기 증거로 해석해야 하며, 사람 평가를 대체하는 최종 증거로 보아서는 안 된다.

\subsection{EmoTrace가 실제로 만든 것}

EmoTrace는 더 긍정적인 응답 생성기를 만든 것이 아니다. 오히려 기존 LLM 응답이 긍정성과 부드러움으로 기울어지는 문제를 드러냈고, 내부 trace와 episode 정보를 사용해 특정 조건에서 그 편향을 줄일 수 있음을 보였다.

따라서 EmoTrace의 산출물은 세 가지로 정리된다. 첫째, 감정을 내부 동역학으로 보내는 실행 가능한 파이프라인이다. 둘째, trace를 감정 상태 표현 후보로 분석하기 위한 지표 체계이다. 셋째, trace를 episode로 해석해 응답 조건으로 사용하는 적용 구조이다. 이 세 요소가 결합되어 EmoTrace는 감정 분류기, 단순 정서 벡터 모델, 프롬프트 기반 공감 응답과 구분된다.

\section{한계와 주장 가능 범위}

\begin{figure}[t]
  \centering
  \includegraphics[width=0.9\linewidth]{\figdir/claim_boundary_map.png}
  \caption{현재 주장 가능 범위와 아직 주장하면 안 되는 범위.}
  \label{fig:academic-claim-boundary}
\end{figure}

현재 주장 가능한 것은 다음과 같다. 첫째, v3 보정 이후 branch collapse가 완화되었다. 둘째, v3.1 trace는 감정축 정보를 일부 담고 조작에 반응했다. 셋째, v4 targeted 조건에서는 episode trace가 자극만 사용하는 기준선보다 에피소드 충실도를 높였다.

반대로 아직 주장하면 안 되는 것도 분명하다. EmoTrace가 인간처럼 실제 감정을 느낀다고 주장할 수 없다. 모든 일반 감정 응답에서 항상 우수하다고 주장할 수 없다. LLM judge 결과만으로 사람 평가를 통과했다고 볼 수 없다. v1/v2의 정량 성능, 뉴런 타입별 ablation, 대규모 사용자 평가도 아직 충분하지 않다.

\section{향후 연구}

향후 연구는 다섯 방향으로 진행되어야 한다. 첫째, 사람 평가를 통해 LLM judge 편향을 검증해야 한다. 둘째, v1/v2의 정량 실험과 GUI 스크린샷을 보강해야 한다. 셋째, 흥분성, 억제성, 조절성 뉴런을 제거하는 ablation 실험을 수행해야 한다. 넷째, trace 인과 검증을 더 큰 규모로 확장해야 한다. 다섯째, 장기 감정 상태와 episode memory를 도입해 반복 상호작용에서 감정 상태가 어떻게 변하는지 분석해야 한다.

\section{결론}

본 연구는 LLM 공감 응답의 어색함에서 출발하여, 감정을 출력 라벨이나 문체가 아니라 내부 신경 활성 trace로 다루는 EmoTrace 구조를 제안하였다. v1은 감정을 내부 상태로 보내는 첫 파이프라인을 만들었고, v2는 모듈형 뉴런망과 branch 기록을 구성하였다. v3는 dominant branch를 추출하면서 branch collapse와 스타일 편향을 발견하였고, v3.1은 감정 정의를 바꾸어 trace 자체를 감정 상태 표현 후보로 보았다. v4는 trace를 episode로 해석하고 targeted 조건에서 응답 충실도 향상을 확인하였다.

최종적으로 EmoTrace는 모든 감정 응답을 해결한 완성형 시스템이 아니다. 그러나 자극을 내부 신경 활성 흐름으로 전개하고, 그 흐름을 감정 상태 표현 후보로 분석하며, targeted episode-sensitive 조건에서 자극만 사용하는 기준선보다 감정 에피소드 충실도를 높일 수 있음을 보였다. 이것이 현재 가장 방어 가능한 결론이다.

\appendix

\section{핵심 용어}

\begin{longtable}{p{0.24\linewidth}p{0.66\linewidth}}
\toprule
용어 & 의미 \\
\midrule
자극 & 입력 문장이 내부 동역학에 가하는 최초 압력. v3.1 이후 초기 감정 벡터는 감정이 아니라 자극으로 해석된다. \\
trace & 자극이 뉴런망 안에서 만든 시간적 활성 기록. 감정 상태 표현 후보이다. \\
지배 흐름 & trace에서 추출한 대표 경로. \\
\eng{z} & 내부 흐름 또는 히스토리를 압축한 잠재 표현. \\
\eng{s} & 응답 스타일 또는 표면화 조건으로 쓰이는 표현. \\
감정 에피소드 & trace를 사건 원인, 평가 구조, 원초 정서, 행동 경향으로 해석한 중간 표현. \\
과도한 온건화 & LLM 응답이 감정의 거친 핵심을 지나치게 부드럽게 덮는 현상. \\
\bottomrule
\end{longtable}

\section{모듈과 증거의 연결}

\begin{figure}[h]
  \centering
  \includegraphics[width=\linewidth]{\figdir/module_evidence_map.png}
  \caption{모듈과 증거의 연결.}
\end{figure}

\section{참고 자료}

\begin{thebibliography}{9}
\bibitem{ekman1992}
Ekman, P. \textit{An Argument for Basic Emotions}. Cognition and Emotion, 6(3--4), 169--200, 1992.

\bibitem{plutchik1980}
Plutchik, R. \textit{A General Psychoevolutionary Theory of Emotion}. In R. Plutchik and H. Kellerman, editors, Theories of Emotion, pp. 3--33. Academic Press, 1980.

\bibitem{russell1980}
Russell, J. A. \textit{A Circumplex Model of Affect}. Journal of Personality and Social Psychology, 39(6), 1161--1178, 1980.

\bibitem{scherer2009}
Scherer, K. R. \textit{The Dynamic Architecture of Emotion: Evidence for the Component Process Model}. Cognition and Emotion, 23(7), 1307--1351, 2009.

\bibitem{gross1998}
Gross, J. J. \textit{Antecedent- and Response-Focused Emotion Regulation: Divergent Consequences for Experience, Expression, and Physiology}. Journal of Personality and Social Psychology, 74(1), 224--237, 1998.

\bibitem{picard1997}
Picard, R. W. \textit{Affective Computing}. MIT Press, 1997.

\bibitem{rashkin2019}
Rashkin, H., Smith, E. M., Li, M., and Boureau, Y.-L. \textit{Towards Empathetic Open-domain Conversation Models: A New Benchmark and Dataset}. ACL, 2019.

\bibitem{zheng2023}
Zheng, L., Chiang, W.-L., Sheng, Y., Zhuang, S., Wu, Z., Zhuang, Y., Lin, Z., Li, Z., Li, D., Xing, E. P., Zhang, H., Gonzalez, J. E., and Stoica, I. \textit{Judging LLM-as-a-Judge with MT-Bench and Chatbot Arena}. NeurIPS Datasets and Benchmarks, 2023.

\bibitem{v2readme}
EmoTrace v2 README. \textit{모듈형 PyTorch MVP 설명 문서}. 2026.

\bibitem{v3draft}
EmoTrace v3 Korean Draft. \textit{감정 동역학 기반 한국어 응답 생성 프레임워크}. 2026.

\bibitem{branchcollapse}
Branch Collapse Mitigation. \textit{v3 branch length collapse 분석 및 완화 기록}. 2026.

\bibitem{v31report}
v3.1 Final Completion Report. \textit{trace-as-emotion 표현 및 인과 검증 보고서}. 2026.

\bibitem{v4summary}
EmoTrace v4 Research Summary. \textit{episode interpretation 및 generation 평가 요약}. 2026.

\bibitem{targeted}
EmoTrace Targeted Superiority Report. \textit{targeted episode-sensitive 입력에서의 쌍대 우월성 평가}. 2026.
\end{thebibliography}

\end{document}


